\section{Introduzione}
\subsection{Parallel n-grams}

\begin{frame}{Introduzione}
        Conteggio degli $n$-grams (sequenze di $n$ parole consecutive) in file di testo.
        \begin{itemize}
            \item Input:
            \begin{itemize}
                \item File di testo inglese (normalizzazione semplificata).
                \item Lunghezza $n$ degli $n$-grams (default 3).
            \end{itemize}
            \item Output: Statistiche sugli $n$-grams più frequenti e sui tempi di esecuzione.
        \end{itemize}
    \begin{block}{Versioni}
        Sono state realizzate due implementazioni:
        \begin{itemize}
            \item Sequenziale (C).
            \item Parallela (OpenMP).
        \end{itemize}
    \end{block}
\end{frame}

\begin{frame}{Introduzione}
    \begin{block}{Input}
        Dataset testuale (es. \href{https://huggingface.co/datasets/allenai/c4/tree/main/en}{Hugging Face C4}).
        \begin{itemize}
            \item Script Python per il download. Permette di specificare la dimensione del file di input.
        \end{itemize}
    \end{block}
    \begin{block}{Normalizzazione}
        Creazione di classi di equivalenza per parole simili:
        \begin{itemize}
            \item Tutto in minuscolo.
            \item Rimozione punteggiatura.
            \item Si mantengono solo caratteri alfanumerici e singolo spazio.
        \end{itemize}
    \end{block}
\end{frame}

\begin{frame}{Introduzione}
    \begin{block}{Struttura Dati}
        Utilizzo di una Hash Table per contare gli $n$-grams.
        \begin{itemize}
            \item Hash Function: Polynomial Hash,
                \begin{equation*}
                    H(c)=\left(\sum_{i=1}^n c_i\cdot p^{n-i}\right) \mod M
                \end{equation*}
            \item Hash Function implementation: Dato $H(c_1\cdots c_i)$,
                \begin{equation*}
                    H(c_1\cdots c_ic_{i+1})=(H(c_1\cdots c_i)\cdot p+c_{i+1}) \mod M.
                \end{equation*}
            \item Collisioni: Open Chain.
        \end{itemize}
    \end{block}
\end{frame}

\begin{frame}{Introduzione}
    \begin{block}{Struttura Dati}
        Hash Function: $\displaystyle H(c)=\left(\sum_{i=1}^n c_i\cdot p^{n-i}\right) \mod M$
        \begin{itemize}
            \item $p$: Deve essere primo e più grande della dimensione dell'alfabeto. Abbiamo scelto $p=257$.
            \item $M$: Deve essere primo. Regoalre $M$ perfette di ottenere un fill factor $\alpha=\tfrac{busy\_buckets}{table\_length}$ migliore o peggiore. Abbiamo scelto $M=10000019$.
            \item Overflow: Problemi se durante il calcolo $H(c_1,\cdots,c_i)$ è più grande del tipo di dato usato per un qualche $i$. Abbiamo usato il tipo più veloce per l'architettura di almeno 32 bit (i.e., \texttt{uint\_fast32\_t}): $inter\_var\_dim=(M-1)\cdot p+c_{max}<2^{32}$.
        \end{itemize}
    \end{block}
\end{frame}
